\color{unidarkgreen}{\subsection{Управление выключателем через УПНКП}}
\color{black}

\begin{enumerate}

\item
Функциональный блок управления выключателем (УВ УПНКП) предназначен для подачи сигнала на включение и отключение силового выключателя ШР через устройство принудительной несинхронной коммутации полюсов (УПНКП). 

\item
Ввод в работу логики управления выключателем через УПНКП осуществляется с помощью программного ключа \textbf{«УВ УПНКП»}.

\item
При наличии сигнала \textit{«КСВ: Блокировка управления: Блок. включения»}, формируемого функциональным блоком контроля силового выключателя, команда на включение выключателя блокируется. При наличии сигнала \textit{«КСВ: Блокировка управления: Блок. отключения»}, формируемого функциональным блоком контроля силового выключателя, команда на отключение выключателя блокируется.

\item
Оперативный ввод/вывод УВ производится с помощью БУ (блоков управления функциями) \textit{«УВ УПНКП: Вывод»}. 

\item
На рисунке \ref{pic:UVupnkp} представлена функциональная схема логики УВ УПНКП. Уставки УВ УПНКП представлены в таблице \ref{app_set:uvupnkp}. 

\medskip
\begin{figure}[H]
\centering
\includegraphics[width=1\textwidth,height=1\textheight,keepaspectratio]{\fbpath/11.01 ШР/6911 УВ УПНКП/img/UV.pdf}
\captionof{figure}{Функциональная схема логики отключения выключателя}\label{pic:UVupnkp}
\end{figure}


\end{enumerate}